\color{unimain}{\subsection{Фиксация положения выключателей (ФПВ) }}
\color{black}

\begin{enumerate}


\item
Логика ФПВ В1(ЛВ), ФПВ В2(ОВ), ФПВ В3, ФПВ В1(ЛВ)СН, ФПВ В2(ОВ)СН аналогичны, рассмотрим на примере ФПВ В1(ЛВ). 

\item
Ввод в работу ФПВ В1(ЛВ)  и  осуществляется с помощью уставок \textbf{«ФПВ В1(ЛВ)»} . Оперативный 
ввод/вывод ФПВ производится с помощью БУ (блоков управления функциями) «Вывод ФПВ В1(ЛВ) ». 
Предусматривается возможность формирования положения выключателей с контролем 
разъединителей, которая задается уставками \textbf{«Контроль Р В1(ЛВ)»}. С помощью уставок \textbf{«Схема Р 
В1(ЛВ)»} имеется возможность выбрать схему выключателя, два или три разъединителя. С помощью 
уставки \textbf{«Тип В1(ЛВ)»} имеется возможность выбрать тип выключателя, линейный или обходной. Так 
же от блок-контактов В1(ЛВ)   формируется  внутренний пуск ЗНФ, который пофазно контролирует 
расхождение контактов полюсов выключателей, возникающее при трехфазном включении или 
отключении выключателя. 


\begin{figure}[H]
\centering
\includegraphics[width=0.75\textwidth,height=0.75\textheight,keepaspectratio]{\fbpath/761.3471 ФПВ/img/FPV_1.pdf}
\end{figure}


\begin{figure}[H]
\centering
\includegraphics[width=0.75\textwidth,height=0.75\textheight,keepaspectratio]{\fbpath/761.3471 ФПВ/img/FPV_2.pdf}
\end{figure}


\begin{figure}[H]
\centering
\includegraphics[width=0.75\textwidth,height=0.75\textheight,keepaspectratio]{\fbpath/761.3471 ФПВ/img/FPV_logic.pdf}
\captionof{figure}{Функциональная схема ФПВ В1}\label{pict_fpv:FPV}
\end{figure}

\end{enumerate} 